% Plantilla reconstruida a partir de "Primer Trabajo Estadistica II.docx".
% Codificacion del archivo: UTF-8.
% Compilador recomendado en Overleaf y Texmaker: pdfLaTeX.
% No requiere shell-escape ni programas externos.

\documentclass[12pt,letterpaper]{article}

% Idioma, codificacion y tipografia compatibles con pdfLaTeX
\usepackage[utf8]{inputenc}
\usepackage[T1]{fontenc}
% "provide=*" permite usar la definicion regional de Babel cuando una
% instalacion local minima no incluye spanish.ldf.
\usepackage[spanish,provide=*]{babel}
\usepackage{lmodern}
\usepackage{microtype}

% Pagina y elementos graficos
\usepackage[margin=2.5cm]{geometry}
\usepackage{graphicx}
\graphicspath{{figuras/}}
\usepackage{float}
\usepackage{caption}
\usepackage[table]{xcolor}

% Tablas, listas y codigo
\usepackage{array}
\usepackage{booktabs}
\usepackage{tabularx}
\usepackage{enumitem}
\usepackage{listings}

% Encabezados, pies y enlaces
\usepackage{fancyhdr}
\usepackage{lastpage}
\usepackage[hidelinks]{hyperref}

\definecolor{azulcurso}{HTML}{24495A}
\definecolor{naranjacurso}{HTML}{D9783D}
\definecolor{grisclaro}{HTML}{EDF2F3}
\definecolor{gristexto}{HTML}{555555}

\hypersetup{
  pdftitle={Primer Trabajo Estadistica II},
  pdfsubject={Inferencia estadistica con datos de la GEIH 2025}
}

\pagestyle{fancy}
\fancyhf{}
\fancyhead[L]{\small Estadistica II}
\fancyhead[R]{\small Primer trabajo}
\fancyfoot[C]{\small Pagina \thepage\ de \pageref{LastPage}}
\setlength{\headheight}{15pt}

\setlength{\parindent}{0pt}
\setlength{\parskip}{6pt}
\renewcommand{\arraystretch}{1.25}

\newcolumntype{Y}{>{\raggedright\arraybackslash}X}
\newcommand{\respuesta}[1][Escriba aqui]{\textcolor{gristexto}{\textit{#1}}}

\lstdefinestyle{codigoR}{
  language=R,
  basicstyle=\ttfamily\small,
  keywordstyle=\color{azulcurso}\bfseries,
  commentstyle=\color{gristexto}\itshape,
  stringstyle=\color{naranjacurso},
  numbers=left,
  numberstyle=\tiny\color{gristexto},
  stepnumber=1,
  numbersep=8pt,
  frame=single,
  rulecolor=\color{grisclaro},
  backgroundcolor=\color{grisclaro!45},
  breaklines=true,
  showstringspaces=false,
  columns=fullflexible,
  keepspaces=true,
  literate=
    {á}{{\'a}}1 {é}{{\'e}}1 {í}{{\'i}}1 {ó}{{\'o}}1 {ú}{{\'u}}1
    {Á}{{\'A}}1 {É}{{\'E}}1 {Í}{{\'I}}1 {Ó}{{\'O}}1 {Ú}{{\'U}}1
    {ñ}{{\~n}}1 {Ñ}{{\~N}}1 {ü}{{\"u}}1 {Ü}{{\"U}}1
}

\newcommand{\figuraejemplo}[3]{%
  \begin{figure}[H]
    \centering
    \IfFileExists{figuras/#1}{%
      \includegraphics[width=#2\textwidth]{#1}%
    }{%
      \fbox{\parbox[c][5cm][c]{0.88\textwidth}{\centering
      Inserte aqui la figura correspondiente.\\[4pt]
      Archivo esperado: \texttt{figuras/#1}}}%
    }
    \caption{#3}
  \end{figure}
}

\newcommand{\titulovariable}[1]{%
  \section{\texorpdfstring{#1}{Variable}}
}

\begin{document}

\begin{center}
  {\LARGE\bfseries\color{azulcurso} Primer Trabajo Estadistica II}\par
  \vspace{3pt}
  {\large Inferencia Estadistica}
\end{center}

\vspace{8pt}
\textbf{Integrantes:}
\begin{enumerate}[leftmargin=1.6cm]
  \item \respuesta[Nombre]
  \item \respuesta[Nombre]
  \item \respuesta[Nombre]
\end{enumerate}

\textbf{Fuente de los datos (URL):} \respuesta

\textbf{Descripcion de los datos seleccionados:}

\respuesta[Describa los meses, archivos, poblacion y unidad de observacion empleados.]

\subsection*{Variables seleccionadas}

\textbf{Variable cualitativa}

\begin{tabularx}{\textwidth}{>{\bfseries}p{1.2cm} Y Y}
\toprule
No. & Nombre de la variable & Categorias o niveles \\
\midrule
1 & \respuesta & \respuesta \\
\bottomrule
\end{tabularx}

\vspace{8pt}
\textbf{Variables cuantitativas}

\begin{tabularx}{\textwidth}{>{\bfseries}p{1.2cm} Y Y}
\toprule
No. & Nombre de la variable & Unidad de medicion \\
\midrule
1 & \respuesta & \respuesta \\
2 & \respuesta & \respuesta \\
3 & \respuesta & \respuesta \\
\bottomrule
\end{tabularx}

\section{Desarrollo}

% -----------------------------------------------------------------------------
% VARIABLE 1
% -----------------------------------------------------------------------------
\subsection{\texorpdfstring{Variable 1: \respuesta[Nombre de la variable]}{Variable 1}}

\subsubsection*{Analisis descriptivo}
\textbf{Descriptivos basicos}

\begin{tabularx}{\textwidth}{>{\bfseries}p{1.2cm} Y Y}
\toprule
No. & Medida & Valor \\
\midrule
1 & \respuesta & \respuesta \\
2 & \respuesta & \respuesta \\
3 & \respuesta & \respuesta \\
4 & \respuesta & \respuesta \\
5 & \respuesta & \respuesta \\
6 & \respuesta & \respuesta \\
7 & \respuesta & \respuesta \\
8 & \respuesta & \respuesta \\
9 & \respuesta & \respuesta \\
10 & \respuesta & \respuesta \\
\bottomrule
\end{tabularx}

\textit{Nota: calcule solamente las medidas apropiadas para el tipo de variable.}

\textbf{Comentarios:} \respuesta[Interprete los resultados en el contexto del problema.]

\subsubsection*{Analisis grafico}
\figuraejemplo{ejemplo_analisis_grafico.png}{0.62}{Ejemplo de representacion grafica. Reemplace esta imagen por la correspondiente a la variable.}

\textbf{Comentarios:} \respuesta[Interprete el grafico.]

\subsubsection*{Calculo de los estimadores}

\begin{tabularx}{\textwidth}{p{2.5cm} Y Y Y Y}
\toprule
\textbf{Estimador} & \multicolumn{2}{c}{\textbf{Estimaciones puntuales}} & \multicolumn{2}{c}{\textbf{Estimacion por intervalo}} \\
\cmidrule(lr){2-3}\cmidrule(lr){4-5}
 & \textbf{Analogia} & \textbf{Maxima verosimilitud} & \textbf{Limite inferior} & \textbf{Limite superior} \\
\midrule
\respuesta[Medida] & \respuesta[Resultado] & \respuesta[Resultado] & \respuesta[Resultado] & \respuesta[Resultado] \\
\textbf{Comentario} & \respuesta & \respuesta & \respuesta & \respuesta \\
\bottomrule
\end{tabularx}

\textit{Nota: repita la tabla para cada estimador.}

\subsubsection*{Evaluacion del estimador}

\textbf{Insesgamiento}

\begin{tabularx}{\textwidth}{YYYY}
\toprule
\textbf{Media} & \textbf{Mediana} & \textbf{Varianza} & \textbf{Desviacion estandar} \\
\midrule
\respuesta[Resultado] & \respuesta[Resultado] & \respuesta[Resultado] & \respuesta[Resultado] \\
\bottomrule
\end{tabularx}

\textbf{Comentario:} \respuesta

\textbf{Consistencia}
\figuraejemplo{ejemplo_consistencia.png}{0.50}{Ejemplo para evaluar la consistencia. Reemplace esta imagen por los resultados de la variable.}

\textbf{Comentario:} \respuesta

\textbf{Eficiencia}

\begin{tabularx}{0.58\textwidth}{Y Y}
\toprule
\textbf{Medida} & \textbf{Valor} \\
\midrule
Media & \respuesta \\
Mediana & \respuesta \\
Varianza & \respuesta \\
Desviacion estandar & \respuesta \\
\bottomrule
\end{tabularx}

\textbf{Comentario:} \respuesta

\subsubsection*{Sintaxis empleada con esta variable}
\begin{lstlisting}[style=codigoR]
# Coloque aqui la sintaxis empleada para procesar esta variable.
\end{lstlisting}

% -----------------------------------------------------------------------------
% VARIABLE 2
% -----------------------------------------------------------------------------
\titulovariable{Variable 2: \respuesta[Nombre de la variable]}

\subsection*{Analisis descriptivo}
\textbf{Descriptivos basicos}

\begin{tabularx}{\textwidth}{>{\bfseries}p{1.2cm} Y Y}
\toprule
No. & Medida & Valor \\
\midrule
1 & \respuesta & \respuesta \\
2 & \respuesta & \respuesta \\
3 & \respuesta & \respuesta \\
4 & \respuesta & \respuesta \\
5 & \respuesta & \respuesta \\
6 & \respuesta & \respuesta \\
7 & \respuesta & \respuesta \\
8 & \respuesta & \respuesta \\
9 & \respuesta & \respuesta \\
10 & \respuesta & \respuesta \\
\bottomrule
\end{tabularx}

\textit{Nota: calcule solamente las medidas apropiadas para el tipo de variable.}

\textbf{Comentarios:} \respuesta

\subsection*{Analisis grafico}
\figuraejemplo{ejemplo_analisis_grafico.png}{0.62}{Ejemplo de representacion grafica. Reemplace esta imagen por la correspondiente a la variable.}

\textbf{Comentarios:} \respuesta

\subsection*{Calculo de los estimadores}

\begin{tabularx}{\textwidth}{p{2.5cm} Y Y Y Y}
\toprule
\textbf{Estimador} & \multicolumn{2}{c}{\textbf{Estimaciones puntuales}} & \multicolumn{2}{c}{\textbf{Estimacion por intervalo}} \\
\cmidrule(lr){2-3}\cmidrule(lr){4-5}
 & \textbf{Analogia} & \textbf{Maxima verosimilitud} & \textbf{Limite inferior} & \textbf{Limite superior} \\
\midrule
\respuesta[Medida] & \respuesta[Resultado] & \respuesta[Resultado] & \respuesta[Resultado] & \respuesta[Resultado] \\
\textbf{Comentario} & \respuesta & \respuesta & \respuesta & \respuesta \\
\bottomrule
\end{tabularx}

\textit{Nota: repita la tabla para cada estimador.}

\subsection*{Evaluacion del estimador}
\textbf{Insesgamiento}

\begin{tabularx}{\textwidth}{YYYY}
\toprule
\textbf{Media} & \textbf{Mediana} & \textbf{Varianza} & \textbf{Desviacion estandar} \\
\midrule
\respuesta[Resultado] & \respuesta[Resultado] & \respuesta[Resultado] & \respuesta[Resultado] \\
\bottomrule
\end{tabularx}

\textbf{Comentario:} \respuesta

\textbf{Consistencia}
\figuraejemplo{ejemplo_consistencia.png}{0.50}{Ejemplo para evaluar la consistencia. Reemplace esta imagen por los resultados de la variable.}

\textbf{Comentario:} \respuesta

\textbf{Eficiencia}

\begin{tabularx}{0.58\textwidth}{Y Y}
\toprule
\textbf{Medida} & \textbf{Valor} \\
\midrule
Media & \respuesta \\
Mediana & \respuesta \\
Varianza & \respuesta \\
Desviacion estandar & \respuesta \\
\bottomrule
\end{tabularx}

\textbf{Comentario:} \respuesta

\subsection*{Sintaxis empleada con esta variable}
\begin{lstlisting}[style=codigoR]
# Coloque aqui la sintaxis empleada para procesar esta variable.
\end{lstlisting}

% -----------------------------------------------------------------------------
% VARIABLE 3
% -----------------------------------------------------------------------------
\titulovariable{Variable 3: \respuesta[Nombre de la variable]}

\subsection*{Analisis descriptivo}
\textbf{Descriptivos basicos}

\begin{tabularx}{\textwidth}{>{\bfseries}p{1.2cm} Y Y}
\toprule
No. & Medida & Valor \\
\midrule
1 & \respuesta & \respuesta \\
2 & \respuesta & \respuesta \\
3 & \respuesta & \respuesta \\
4 & \respuesta & \respuesta \\
5 & \respuesta & \respuesta \\
6 & \respuesta & \respuesta \\
7 & \respuesta & \respuesta \\
8 & \respuesta & \respuesta \\
9 & \respuesta & \respuesta \\
10 & \respuesta & \respuesta \\
\bottomrule
\end{tabularx}

\textit{Nota: calcule solamente las medidas apropiadas para el tipo de variable.}

\textbf{Comentarios:} \respuesta

\subsection*{Analisis grafico}
\figuraejemplo{ejemplo_analisis_grafico.png}{0.62}{Ejemplo de representacion grafica. Reemplace esta imagen por la correspondiente a la variable.}

\textbf{Comentarios:} \respuesta

\subsection*{Calculo de los estimadores}

\begin{tabularx}{\textwidth}{p{2.5cm} Y Y Y Y}
\toprule
\textbf{Estimador} & \multicolumn{2}{c}{\textbf{Estimaciones puntuales}} & \multicolumn{2}{c}{\textbf{Estimacion por intervalo}} \\
\cmidrule(lr){2-3}\cmidrule(lr){4-5}
 & \textbf{Analogia} & \textbf{Maxima verosimilitud} & \textbf{Limite inferior} & \textbf{Limite superior} \\
\midrule
\respuesta[Medida] & \respuesta[Resultado] & \respuesta[Resultado] & \respuesta[Resultado] & \respuesta[Resultado] \\
\textbf{Comentario} & \respuesta & \respuesta & \respuesta & \respuesta \\
\bottomrule
\end{tabularx}

\textit{Nota: repita la tabla para cada estimador.}

\subsection*{Evaluacion del estimador}
\textbf{Insesgamiento}

\begin{tabularx}{\textwidth}{YYYY}
\toprule
\textbf{Media} & \textbf{Mediana} & \textbf{Varianza} & \textbf{Desviacion estandar} \\
\midrule
\respuesta[Resultado] & \respuesta[Resultado] & \respuesta[Resultado] & \respuesta[Resultado] \\
\bottomrule
\end{tabularx}

\textbf{Comentario:} \respuesta

\textbf{Consistencia}
\figuraejemplo{ejemplo_consistencia.png}{0.50}{Ejemplo para evaluar la consistencia. Reemplace esta imagen por los resultados de la variable.}

\textbf{Comentario:} \respuesta

\textbf{Eficiencia}

\begin{tabularx}{0.58\textwidth}{Y Y}
\toprule
\textbf{Medida} & \textbf{Valor} \\
\midrule
Media & \respuesta \\
Mediana & \respuesta \\
Varianza & \respuesta \\
Desviacion estandar & \respuesta \\
\bottomrule
\end{tabularx}

\textbf{Comentario:} \respuesta

\subsection*{Sintaxis empleada con esta variable}
\begin{lstlisting}[style=codigoR]
# Coloque aqui la sintaxis empleada para procesar esta variable.
\end{lstlisting}

% -----------------------------------------------------------------------------
% VARIABLE 4
% -----------------------------------------------------------------------------
\titulovariable{Variable 4: \respuesta[Nombre de la variable]}

\subsection*{Analisis descriptivo}
\textbf{Descriptivos basicos}

\begin{tabularx}{\textwidth}{>{\bfseries}p{1.2cm} Y Y}
\toprule
No. & Medida & Valor \\
\midrule
1 & \respuesta & \respuesta \\
2 & \respuesta & \respuesta \\
3 & \respuesta & \respuesta \\
4 & \respuesta & \respuesta \\
5 & \respuesta & \respuesta \\
6 & \respuesta & \respuesta \\
7 & \respuesta & \respuesta \\
8 & \respuesta & \respuesta \\
9 & \respuesta & \respuesta \\
10 & \respuesta & \respuesta \\
\bottomrule
\end{tabularx}

\textit{Nota: calcule solamente las medidas apropiadas para el tipo de variable.}

\textbf{Comentarios:} \respuesta

\subsection*{Analisis grafico}
\figuraejemplo{ejemplo_analisis_grafico.png}{0.62}{Ejemplo de representacion grafica. Reemplace esta imagen por la correspondiente a la variable.}

\textbf{Comentarios:} \respuesta

\subsection*{Calculo de los estimadores}

\begin{tabularx}{\textwidth}{p{2.5cm} Y Y Y Y}
\toprule
\textbf{Estimador} & \multicolumn{2}{c}{\textbf{Estimaciones puntuales}} & \multicolumn{2}{c}{\textbf{Estimacion por intervalo}} \\
\cmidrule(lr){2-3}\cmidrule(lr){4-5}
 & \textbf{Analogia} & \textbf{Maxima verosimilitud} & \textbf{Limite inferior} & \textbf{Limite superior} \\
\midrule
\respuesta[Medida] & \respuesta[Resultado] & \respuesta[Resultado] & \respuesta[Resultado] & \respuesta[Resultado] \\
\textbf{Comentario} & \respuesta & \respuesta & \respuesta & \respuesta \\
\bottomrule
\end{tabularx}

\textit{Nota: repita la tabla para cada estimador.}

\subsection*{Evaluacion del estimador}
\textbf{Insesgamiento}

\begin{tabularx}{\textwidth}{YYYY}
\toprule
\textbf{Media} & \textbf{Mediana} & \textbf{Varianza} & \textbf{Desviacion estandar} \\
\midrule
\respuesta[Resultado] & \respuesta[Resultado] & \respuesta[Resultado] & \respuesta[Resultado] \\
\bottomrule
\end{tabularx}

\textbf{Comentario:} \respuesta

\textbf{Consistencia}
\figuraejemplo{ejemplo_consistencia.png}{0.50}{Ejemplo para evaluar la consistencia. Reemplace esta imagen por los resultados de la variable.}

\textbf{Comentario:} \respuesta

\textbf{Eficiencia}

\begin{tabularx}{0.58\textwidth}{Y Y}
\toprule
\textbf{Medida} & \textbf{Valor} \\
\midrule
Media & \respuesta \\
Mediana & \respuesta \\
Varianza & \respuesta \\
Desviacion estandar & \respuesta \\
\bottomrule
\end{tabularx}

\textbf{Comentario:} \respuesta

\subsection*{Sintaxis empleada con esta variable}
\begin{lstlisting}[style=codigoR]
# Coloque aqui la sintaxis empleada para procesar esta variable.
\end{lstlisting}

\end{document}
